\cotefile{57}
\scanpage{242}\mspage{- 16 -}
n$^{\mathrm{o}}$ 227

3. \uline{L'application rationnelle définie par un système linéaire}.

On continue à supposer $V$ complète. Soit $\handwritten{\mathscr L}$ un système linéaire,
$X_{\unsure{0|o}}\handwritten{\in\mathscr L}$, et $L_0$ le module dérivé de $\handwritten{\mathscr L}$ par $X_0$. Soit $f_0,\ldots,f_n$ une
base de ce module. Alors les fonctions $(\handwritten{1},f_{\handwritten{1}}/f_0\unsure{,|;}\ldots,f_n/f_0)$ peuvent
être prises comme coordonnées non homogènes d'un point dans l'espace
projectif, qui représente un point générique d'une variété projective
$V'$. Si $t$ est un élément transcendant sur tout, alors $(tf_0,\ldots,tf_n)$
sont les coordonnées homogènes d'un point générique de $V'$. Elles
définissent une application rationnelle $T:V\handwritten{\longrightarrow}V'$ de $V$ sur $V'$.
Un changement de base dans $L_0$ donne une transformation projective de $V'$.

Si on employait un module $L_{\handwritten{1}}$ dérivé de $\handwritten{\mathscr L}$ par un autre diviseur
$X_{\handwritten{1}}$, on voit que l'application rationnelle $T_{\handwritten{1}}$ obtenue comme ci-dessus
diffère de $T$ par une transformation \deleted{linéaire}\inserted{\handwritten{\unsure{projective}}}. On peut donc dire que $T$
est uniquement déterminée par $\handwritten{\mathscr L}$. De plus, si $\handwritten{\mathscr L}$ et $\handwritten{\mathscr L'}$ sont deux
\margin{\handwritten{\(\big|\)}}systèmes linéaires qui ne diffèrent que par une composante fixe, alors
on voit immédiatement qu'ils donnent lieu à la même application rationnelle
$T$ de $V$, les modules dérivés étant les mêmes. Pour étudier l'application
rationnelle associée à un système linéaire, on peut donc se
borner, sans perte de généralité, aux systèmes linéaires sans composantes
fixes.

On va maintenant donner un critère pour que l'application rationnelle
$T$ soit définie en un point $P$ de $V$.

Un point $P$ de $V$ est dit un \uline{point de base} du système linéaire $\handwritten{\mathscr L}$
si $P\handwritten{\in\unsure{\gamma|r}}\handwritten{|}X\handwritten{|}$ pour tout $X\handwritten{\in\mathscr L}$. Si $Y$ est une composante fixe de $\handwritten{\mathscr L}$, alors
tous les points de $\handwritten{|}Y\handwritten{|}$ sont des points de base. Il peut très bien
arriver qu'un système linéaire n'ait pas de composante fixe, mais ait
% archive sheet number: 241
